% ATENÇÃO - veja com o seu orientador se você vai ter este capítulo e se este vai ter nome!
\chapter{Trabalhos Relacionados}
\label{cap:trabalhos:relacionados}

Diversos trabalhos avaliam o gRPC por ângulos diferentes. No campo do desempenho,
\citeonline{PerformanceEvaluation:2024} compararam REST, GraphQL e gRPC em microsserviços em Go,
sob 100 a 500 requisições concorrentes, e o gRPC obteve os menores tempos de resposta e o menor uso
de CPU. \citeonline{Keeps-CaseStudy:2025} relata um caso em que a migração das chamadas síncronas
internas para gRPC, combinada com cache e outras melhorias, reduziu a latência p50 em
cerca de 50\% e aumentou a vazão em 114\%. Ambos avaliam apenas desempenho, e o segundo não isola o
efeito do gRPC.

Outros trabalhos examinam o custo de intermediários na comunicação.
\citeonline{LoadBalancer-Evaluation:2022} mostrou que um proxy L7 no balanceamento de
carga gRPC adiciona cerca de 1 ms por chamada e consome mais CPU e memória do que o balanceamento
no cliente. \citeonline{RethinkingRPC:2025} apontam o empilhamento de protocolos e os
\textit{sidecars} como fontes de latência e consumo de CPU, e propõem uma pilha RPC compilada, com
descarregamento de funções de rede para eBPF e hardware programável. Esse custo é relevante aqui,
pois o Envoy também intermedeia as chamadas, mas nenhum dos dois avalia os mecanismos do próprio gRPC.

Por fim, há trabalhos voltados à resiliência. \citeonline{ResilienceBench-Operator:2025} propõem
uma extensão do Kubernetes que automatiza a avaliação de estratégias de resiliência, com injeção de
falhas via Envoy em chamadas HTTP e gRPC. \citeonline{FastFI:2026} propõem o FastFI, que
usa injeção de falhas combinatórias para decidir em quais chamadas aplicar \textit{retries} e
\textit{fallbacks}, e reduz em mais de 76\% o tempo dos experimentos em relação às soluções
existentes. Os dois focam nas ferramentas de teste, e não na avaliação de um sistema completo.

Em resumo, os trabalhos revisados tratam cada um de um aspecto isolado: o desempenho do gRPC, o
custo de um proxy ou a resiliência sob falhas. Entre eles, nenhum avalia sistematicamente
os mecanismos nativos do gRPC (\textit{deadlines}, \textit{retries}, \textit{hedging} e
\textit{keepalive}) sob carga e sob falha, nem seus efeitos conjuntos sobre desempenho, tolerância
a falhas e segurança. É nessa lacuna que este trabalho se posiciona.


%---------------------------------------------------%
